% Bloc — COMP204 Distributed Systems Poster
\documentclass[final]{beamer}

% ====================
% Packages
% ====================

\usepackage[T1]{fontenc}
\usepackage[utf8]{luainputenc}
\usepackage[defaultsans]{opensans}
\renewcommand{\familydefault}{\sfdefault}
\usepackage[size=custom, width=122,height=91, scale=0.95]{beamerposter}
\usetheme{gemini}
\usecolortheme{gemini}
\usepackage{graphicx}
\usepackage{booktabs}
\usepackage{tikz}
\usepackage{pgfplots}
\pgfplotsset{compat=1.14}
\usepackage{anyfontsize}
\usepackage{listings}
\usepackage{xcolor}
\usetikzlibrary{arrows.meta, positioning, shapes.geometric}

% Code listing style — JavaScript
\lstset{
    language=Java,
    basicstyle=\ttfamily\footnotesize,
    keywordstyle=\color{blue},
    commentstyle=\color{green!60!black},
    stringstyle=\color{red},
    breaklines=true,
    frame=single,
    numbers=left,
    numberstyle=\tiny\color{gray},
    morekeywords={const,let,async,await,req,res,next,jwt,bcrypt}
}

% ====================
% Lengths
% ====================

\newlength{\sepwidth}
\newlength{\colwidth}
\setlength{\sepwidth}{0.025\paperwidth}
\setlength{\colwidth}{0.3\paperwidth}

\newcommand{\separatorcolumn}{\begin{column}{\sepwidth}\end{column}}

% ====================
% Title
% ====================

\title{Bloc: A Distributed Bouldering Session Tracker with Dual-Client Architecture}

\author{Alfie James Jenkins}

\institute[shortinst]{Falmouth University -- COMP204 Distributed Systems}

% ====================
% Footer
% ====================

\footercontent{
  \href{https://github.falmouth.ac.uk/AJ313798/bouldering-blog}{github.falmouth.ac.uk/AJ313798/bouldering-blog} \hfill
  COMP204 Distributed Systems -- 2025/26 \hfill
  \href{mailto:aj313798@falmouth.ac.uk}{aj313798@falmouth.ac.uk}}

% ====================
% Body
% ====================

\begin{document}

\begin{frame}[t,fragile]
\begin{columns}[t]
\separatorcolumn

% ============ COLUMN 1 ============
\begin{column}{\colwidth}

  \begin{block}{The Bloc System}
    Bloc tracks bouldering sessions. Climbers log gym visits and individual
    climbs from a browser or a phone; all data goes to the same backend, so
    both views stay in sync.

    \vspace{0.15cm}
    The project is built around one distributed systems pattern:
    \textbf{multiple independent clients sharing a common server and data store}.
    Neither client holds state locally --- both read and write through the same
    REST API, so a session started on the web shows up immediately in the
    mobile app.

    \vspace{0.15cm}
    \textbf{Hosted on:} a Linux VPS running Node.js
  \end{block}

  \vspace{-0.7cm}
  \begin{block}{System Architecture}

    \begin{center}
    \begin{tikzpicture}[
      font=\small,
      box/.style={draw, rounded corners, fill=gray!12, inner sep=10pt,
                  text width=3.4cm, align=center},
      sbox/.style={draw, rounded corners, fill=blue!12, inner sep=10pt,
                   text width=3.4cm, align=center},
      dbbox/.style={draw, cylinder, shape border rotate=90, fill=orange!15,
                    inner sep=8pt, align=center, aspect=0.3},
      arr/.style={latex-latex, thick},
    ]
      \node[box]  (web) at (0, 1.6)  {Web Browser\\[2pt]{\footnotesize Handlebars\\cookie session}};
      \node[box]  (mob) at (0, -1.6) {Mobile App\\[2pt]{\footnotesize React Native\\JWT bearer}};
      \node[sbox] (api) at (5.5, 0)  {Express Server\\[2pt]{\footnotesize Node.js\\REST API}};
      \node[dbbox](db)  at (10.5, 0) {SQLite\\DB};

      \draw[arr] (web.east) -- node[above,sloped,font=\footnotesize]{HTTP/HTTPS} (api.north west);
      \draw[arr] (mob.east) -- node[below,sloped,font=\footnotesize]{HTTPS + JWT} (api.south west);
      \draw[arr] (api.east) -- node[above,font=\footnotesize]{SQL queries} (db.west);
    \end{tikzpicture}
    \end{center}

    \vspace{0.1cm}
    Browsers handle cookies automatically, so the web app uses cookie-based
    sessions. The mobile app has no shared cookie jar --- it sends a JWT in
    every \texttt{Authorization: Bearer} header instead. Both clients hit the
    same database rows.
  \end{block}

  \begin{block}{Technology Stack}
    \begin{center}
    \begin{tabular}{ll}
    \toprule
    \textbf{Layer} & \textbf{Technology} \\
    \midrule
    Mobile client   & React Native, Expo, TypeScript \\
    Web client      & Express, Handlebars, CSS \\
    REST API        & Express (Node.js) \\
    Database        & SQLite3 \\
    Auth (web)      & bcrypt + express-session \\
    Auth (mobile)   & bcrypt + JWT (jsonwebtoken) \\
    File uploads    & Multer (5\,MB limit, image validation) \\
    Hosting         & Linux VPS (Node.js, port 3000) \\
    \bottomrule
    \end{tabular}
    \end{center}
  \end{block}

  \begin{block}{Core Features}
    \begin{itemize}\setlength\itemsep{2pt}
      \item \textbf{Timed sessions} --- start a session at a named gym;
            \texttt{end\_time} is stamped when the user finishes
      \item \textbf{Climb logging} --- grade (V-scale), attempts, topped
            status, zones reached, optional photo per climb
      \item \textbf{Photo uploads} --- \texttt{multipart/form-data} via
            Multer; stored in \texttt{public/uploads/} and served as
            static files; nullable in the schema
      \item \textbf{Public feed} --- all completed sessions visible to
            every user; lets climbers see what others at the gym have logged
      \item \textbf{Shared state} --- web and mobile read the same rows;
            no synchronisation layer required
      \item \textbf{Single active session} --- enforced at query level
            (\texttt{end\_time IS NULL}), consistent across both clients
    \end{itemize}
  \end{block}

  \begin{block}{Future Development}
    \begin{itemize}\setlength\itemsep{2pt}
      \item \textbf{Grade progression charts} --- visualise topped percentage
            per grade over time using session history already in the database
      \item \textbf{Object storage for images} --- replace \texttt{public/uploads/}
            with S3-compatible storage to decouple files from the server process
      \item \textbf{PostgreSQL migration} --- \texttt{db.js} is the only module
            to change; the rest of the stack is database-agnostic
      \item \textbf{Push notifications} --- alert a user when a friend logs a
            session at the same gym, using the public feed as the data source
      \item \textbf{Offline logging} --- queue climbs locally when the device
            has no signal and sync when connectivity returns
    \end{itemize}
  \end{block}

  \begin{block}{References}
    \footnotesize
    \begin{itemize}\setlength\itemsep{0pt}
      \item Fielding, R.\ T.\ (2000). \textit{Architectural Styles and the Design of Network-based Software Architectures}. University of California, Irvine.
      \item Jones, M.\ B., Bradley, J., \& Sakimura, N.\ (2015). \textit{JSON Web Token (JWT)}. RFC 7519. IETF.
      \item Provos, N., \& Mazieres, D.\ (1999). \textit{A Future-Adaptable Password Scheme}. USENIX Annual Technical Conference.
      \item Express.js documentation: \texttt{expressjs.com}
      \item Expo documentation: \texttt{docs.expo.dev}
      \item React Native documentation: \texttt{reactnative.dev}
      \item Multer documentation: \texttt{github.com/expressjs/multer}
    \end{itemize}
  \end{block}

\end{column}

\separatorcolumn

% ============ COLUMN 2 ============
\begin{column}{\colwidth}

  \begin{block}{REST API Design}
    The mobile API has eight endpoints. All session and climb routes go
    through \texttt{requireApiAuth}, which verifies the JWT and attaches
    \texttt{req.userId} before the route handler runs.

    \vspace{0.1cm}
    \begin{center}
    \begin{tabular}{lll}
    \toprule
    \textbf{Method} & \textbf{Endpoint} & \textbf{Auth} \\
    \midrule
    POST & \texttt{/api/auth/login}        & Public \\
    POST & \texttt{/api/auth/register}     & Public \\
    GET  & \texttt{/api/feed}              & Public \\
    GET  & \texttt{/api/sessions}          & JWT \\
    POST & \texttt{/api/sessions}          & JWT \\
    GET  & \texttt{/api/sessions/:id}      & JWT \\
    POST & \texttt{/api/sessions/:id/end}  & JWT \\
    POST & \texttt{/api/sessions/:id/climbs} & JWT \\
    \bottomrule
    \end{tabular}
    \end{center}

    \vspace{0.1cm}
    \texttt{GET /api/feed} returns the 20 most recent completed sessions with
    climb counts, using a \texttt{LEFT JOIN} so sessions with no climbs still
    appear.
  \end{block}

  \begin{alertblock}{The Mobile Auth Challenge}
    Cookie sessions work for browsers because the browser attaches the cookie
    to every same-origin request automatically. Mobile HTTP clients have no
    shared cookie jar --- each request arrives with no identity attached.

    \vspace{0.1cm}
    Sending the session ID manually would mimic cookie behaviour, but without
    the browser's built-in protections. A JWT with its own expiry is the right
    call for a stateless HTTP client.
  \end{alertblock}

  \begin{exampleblock}{Solution: JWT Authentication Middleware}

\begin{lstlisting}
const requireApiAuth = (req, res, next) => {
  const authHeader = req.headers['authorization'];
  // Header format: "Bearer <token>"
  const token = authHeader && authHeader.split(' ')[1];

  if (!token)
    return res.status(401).json({ error: 'Auth required' });

  jwt.verify(token, JWT_SECRET, (err, decoded) => {
    if (err)
      return res.status(401).json({ error: 'Invalid token' });
    req.userId = decoded.userId; // injected for route handlers
    next();
  });
};
\end{lstlisting}

    \vspace{0.1cm}
    Tokens are signed with \texttt{HS256} and expire after 7 days. The app
    keeps the token in React Context and attaches it to every request; the JWT
    encodes the user identity directly, so no session state is needed on the
    server.
  \end{exampleblock}

  \begin{block}{Database Schema}
    Three tables with cascading foreign keys. All state lives in SQLite;
    neither client caches data locally.

    \vspace{0.2cm}
    \begin{center}
    \begin{tikzpicture}[
      entity/.style={draw, fill=gray!10, rounded corners,
                     inner sep=8pt, font=\footnotesize, align=left},
      rel/.style={-{Latex[length=2mm]}, thick},
    ]
      \node[entity] (users) {
        \textbf{users}\\
        \texttt{id} \ PK\\
        \texttt{email} \ UNIQUE\\
        \texttt{password} \ hashed\\
        \texttt{username}\\
        \texttt{created\_at}
      };

      \node[entity, right=1.8cm of users] (sessions) {
        \textbf{sessions}\\
        \texttt{id} \ PK\\
        \texttt{user\_id} \ FK\\
        \texttt{gym\_name}\\
        \texttt{start\_time}\\
        \texttt{end\_time} \ nullable\\
        \texttt{notes} \ nullable
      };

      \node[entity, right=1.8cm of sessions] (climbs) {
        \textbf{climbs}\\
        \texttt{id} \ PK\\
        \texttt{session\_id} \ FK\\
        \texttt{grade}\\
        \texttt{attempts}\\
        \texttt{topped} \ 0/1\\
        \texttt{zones}\\
        \texttt{image\_path} \ nullable
      };

      \draw[rel] (sessions.west) -- node[above,font=\tiny]{1\hspace{6pt}} (users.east)
                 node[near start, above, font=\tiny]{\hspace{-6pt}N};
      \draw[rel] (climbs.west) -- node[above,font=\tiny]{1\hspace{6pt}} (sessions.east)
                 node[near start, above, font=\tiny]{\hspace{-6pt}N};
    \end{tikzpicture}
    \end{center}

    \vspace{0.1cm}
    \texttt{end\_time IS NULL} means the session is still running. Duration is
    calculated from \texttt{end\_time~-~start\_time} at query time and never
    stored, so editing a session cannot produce stale totals.
  \end{block}

  \begin{block}{Web Client vs Mobile Client}
    \begin{center}
    \begin{tabular}{lll}
    \toprule
    \textbf{Property} & \textbf{Web Browser} & \textbf{Mobile App} \\
    \midrule
    Technology      & Express + Handlebars & React Native / Expo \\
    Auth mechanism  & Cookie session       & JWT Bearer token \\
    Auth middleware & \texttt{requireAuth} & \texttt{requireApiAuth} \\
    Image upload    & HTML form (POST)     & \texttt{FormData} (fetch) \\
    \bottomrule
    \end{tabular}
    \end{center}
    \vspace{0.1cm}
    Both clients share the same database --- the difference is in how they
    authenticate, not what they store.
  \end{block}

\end{column}

\separatorcolumn

% ============ COLUMN 3 ============
\begin{column}{\colwidth}

  \begin{block}{Authentication Sequence}
    Mobile login flow, from tap to authenticated state:

    \vspace{0.1cm}
    \begin{center}
    \resizebox{\linewidth}{!}{%
    \begin{tikzpicture}[
      font=\scriptsize,
      life/.style={draw, thick},
      msg/.style={-{Latex[length=3mm]}, thick},
      ret/.style={-{Latex[length=3mm]}, thick, dashed},
      self/.style={thick},
    ]
      \node (A) at (0,0)    {\textbf{Mobile App}};
      \node (B) at (6.0,0)  {\textbf{Express}};
      \node (C) at (11.5,0) {\textbf{SQLite}};

      \draw[life] (0,-0.35)    -- (0,-9.2);
      \draw[life] (6.0,-0.35)  -- (6.0,-9.2);
      \draw[life] (11.5,-0.35) -- (11.5,-9.2);

      % 1. Login request
      \draw[msg] (0,-1.0) -- node[above]{POST /api/auth/login} (6.0,-1.0);

      % 2. DB query
      \draw[msg] (6.0,-2.0) -- node[above]{SELECT * FROM users WHERE email=?} (11.5,-2.0);

      % 3. DB response
      \draw[ret] (11.5,-2.9) -- node[above]{user row} (6.0,-2.9);

      % 4. bcrypt self-loop
      \draw[self] (6.0,-3.6) -- (6.5,-3.6) -- (6.5,-4.2) -- (6.0,-4.2);
      \draw[msg,self] (6.5,-4.2) -- (6.0,-4.2);
      \node[align=left] at (9.0,-3.9) {bcrypt.compare(pwd, hash)};

      % 5. jwt.sign self-loop
      \draw[self] (6.0,-4.9) -- (6.5,-4.9) -- (6.5,-5.5) -- (6.0,-5.5);
      \draw[msg,self] (6.5,-5.5) -- (6.0,-5.5);
      \node[align=left] at (8.8,-5.2) {jwt.sign(\{userId\}, 7d)};

      % 6. Return token
      \draw[ret] (6.0,-6.3) -- node[above]{\{token, user\}} (0,-6.3);

      % 7. Store token
      \draw[self] (0,-7.1) -- (-0.55,-7.1) -- (-0.55,-7.7) -- (0,-7.7);
      \draw[msg,self] (-0.55,-7.7) -- (0,-7.7);
      \node[align=left] at (2.5,-7.4) {store token in Context};

      % 8. Navigate
      \draw[self] (0,-8.3) -- (-0.55,-8.3) -- (-0.55,-8.9) -- (0,-8.9);
      \draw[msg,self] (-0.55,-8.9) -- (0,-8.9);
      \node[align=left] at (2.1,-8.6) {navigate to home};

    \end{tikzpicture}}% end resizebox
    \end{center}
  \end{block}

  \begin{block}{Key Design Decisions}
    \heading{Session-first workflow}
    You have to start a session before logging a climb. One gym visit, one
    session; climbs cannot exist without one.

    \vspace{0.1cm}
    \heading{One active session per user}
    Before inserting a new session the server checks for an existing row where
    \texttt{end\_time IS NULL}. If one exists, its ID comes back in the
    response so the client goes straight to it.

    \vspace{0.1cm}
    \heading{Dual auth}
    Browsers get a cookie session handled automatically; mobile clients get a
    signed JWT stored in React Context and sent as \texttt{Authorization: Bearer}.
    Same \texttt{users} table, same bcrypt hashes --- only the transport differs.
  \end{block}

  \begin{block}{Log-Climb Sequence (Distributed Flow)}
    Adding a climb crosses three service boundaries in one action: JWT
    auth, file storage, and a database write.

    \vspace{0.1cm}
    \begin{center}
    \resizebox{\linewidth}{!}{%
    \begin{tikzpicture}[
      font=\scriptsize,
      life/.style={draw, thick},
      msg/.style={-{Latex[length=3mm]}, thick},
      ret/.style={-{Latex[length=3mm]}, thick, dashed},
      self/.style={thick},
    ]
      \node (A) at (0,0)    {\textbf{Mobile}};
      \node (B) at (4.5,0)  {\textbf{Server}};
      \node (C) at (8.5,0)  {\textbf{Multer}};
      \node (D) at (13.0,0) {\textbf{SQLite}};

      \draw[life] (0,-0.35)    -- (0,-7.8);
      \draw[life] (4.5,-0.35)  -- (4.5,-7.8);
      \draw[life] (8.5,-0.35)  -- (8.5,-7.8);
      \draw[life] (13.0,-0.35) -- (13.0,-7.8);

      % 1. POST with JWT + multipart
      \draw[msg] (0,-0.9) -- node[above]{POST /api/sessions/:id/climbs + Bearer JWT} (4.5,-0.9);

      % 2. jwt.verify
      \draw[self] (4.5,-1.6) -- (5.0,-1.6) -- (5.0,-2.1) -- (4.5,-2.1);
      \draw[msg,self] (5.0,-2.1) -- (4.5,-2.1);
      \node[align=left] at (7.0,-1.85) {jwt.verify $\to$ req.userId};

      % 3. Multer validate + save
      \draw[msg] (4.5,-2.7) -- node[above]{validate + save image} (8.5,-2.7);

      % 4. Write file
      \draw[self] (8.5,-3.3) -- (9.0,-3.3) -- (9.0,-3.8) -- (8.5,-3.8);
      \draw[msg,self] (9.0,-3.8) -- (8.5,-3.8);
      \node[align=left] at (11.0,-3.55) {write to /uploads/};

      % 5. Return path
      \draw[ret] (8.5,-4.4) -- node[above]{image\_path} (4.5,-4.4);

      % 6. INSERT climb
      \draw[msg] (4.5,-5.0) -- node[above]{INSERT INTO climbs} (13.0,-5.0);

      % 7. climbId
      \draw[ret] (13.0,-5.7) -- node[above]{climbId} (4.5,-5.7);

      % 8. 200 response
      \draw[ret] (4.5,-6.4) -- node[above]{200 \{climbId\}} (0,-6.4);

      % 9. Re-render
      \draw[self] (0,-7.1) -- (-0.5,-7.1) -- (-0.5,-7.7) -- (0,-7.7);
      \draw[msg,self] (-0.5,-7.7) -- (0,-7.7);
      \node[align=left] at (1.8,-7.4) {re-render climb list};

    \end{tikzpicture}}% end resizebox
    \end{center}
  \end{block}


\end{column}

\separatorcolumn
\end{columns}
\end{frame}

\end{document}
